%% Supervised Machine Learning (Classification) -- Mini Project
%% Two-page IEEE conference report.
%%
%% Build:  cd submission/report
%%         tectonic -X compile ieee-report.tex --outdir /tmp/ieeebuild
%% Overleaf also works: upload this folder, IEEEtran.cls ships with Overleaf.
%%
%% If it spills past two pages, drop Table II first, then shorten Section VII.

\documentclass[conference]{IEEEtran}
\IEEEoverridecommandlockouts

\usepackage{cite}
\usepackage{amsmath}
% Pinned explicitly: without it the engine falls back to Computer Modern, which is wider
% than IEEE's Times and silently changes the page count between toolchains.
\usepackage{newtxtext,newtxmath}
\usepackage{graphicx}
\usepackage{booktabs}
\usepackage{textcomp}
\usepackage{xcolor}
\usepackage[hidelinks]{hyperref}

% Two-page limit: default float separation costs roughly a line per float edge.
\setlength{\textfloatsep}{6pt plus 2pt minus 2pt}
\setlength{\floatsep}{6pt plus 2pt minus 2pt}
\setlength{\intextsep}{6pt plus 2pt minus 2pt}
\setlength{\abovecaptionskip}{3pt}
\setlength{\belowcaptionskip}{2pt}
\renewcommand{\arraystretch}{0.95}
\setlength{\abovedisplayskip}{2pt}
\setlength{\belowdisplayskip}{2pt}
\setlength{\abovedisplayshortskip}{1pt}
\setlength{\belowdisplayshortskip}{1pt}

\begin{document}

\title{Beyond the Headline: Calibrated News Article Classifier}

\author{\IEEEauthorblockN{Mohamed Riaz F}
\IEEEauthorblockA{\textit{PES1PGE25DS037} \\
\textit{Department of Data Science and Artificial Intelligence} \\
\textit{PES University}, Bengaluru, India}
}

\maketitle

\begin{abstract}
News topic classifiers usually see only a headline, since that is all an RSS feed
carries. I scraped the full article bodies and measured whether the extra text pays. On
8,001 human-labelled articles across 13 topics, the body raises macro-F1 from 0.712 to
0.771, at McNemar $p = 7.5\times10^{-6}$. The taught families were measured first:
logistic regression 0.752, random forest 0.730, and XGBoost 0.735 at five times the fit
cost, while a majority vote of ten models gained 0.001. A linear support vector machine
over word and character n-grams reached 0.771 and was kept. Once calibrated it declines
19.3\% of test articles and raises accuracy on the rest from 77.7\% to 83.4\%. Test
macro-F1 is 0.751 [0.719, 0.780].\end{abstract}
\begin{IEEEkeywords}
text classification, macro-F1, class imbalance, probability calibration, classification
with rejection, support vector machine
\end{IEEEkeywords}

\section{Introduction and Problem}

A reader who follows \textit{The Hindu}, \textit{The Indian Express} and
\textit{Deccan Herald} sees the same PTI wire story three times, filed under a different
section name in each paper. Merging those feeds needs every article re-filed under one
shared set of 13 topics.

An RSS item carries only a headline, so that is the usual input. A collector that follows
the link also has the body, about 100 times more text. More text need not help: bodies
carry house style, navigation furniture and author biographies, which a classifier can
learn instead of the topic. Whether the body pays is the main question here. The classes
are imbalanced at about $6.7\!:\!1$, from 213 validation articles for politics down to 29
for education, so macro-F1 is the primary metric.
\section{Related Work}

The 13 labels are a flattened subset of IPTC Media Topics~\cite{iptc}. Linear classifiers
over term-weighted features have been the reference method for topic classification since
Joachims~\cite{joachims}. Zadrozny and Elkan~\cite{zadrozny} gave the isotonic method used
here, Guo \textit{et al.}~\cite{guo} the calibration error that scores it, and
Chow~\cite{chow} the rejection option. Dietterich~\cite{dietterich} recommends a paired
test, and Efron and Tibshirani~\cite{efron} the resampling used for the intervals.
\section{Data and Preparation}

A Go service polls 97 RSS and Atom feeds into MongoDB and scrapes the linked bodies. The
corpus holds 14{,}189 articles from 40 publishers, $92.6\%$ with a body, and 8{,}001 carry
a single human label over 13 topics. Per-publisher line counting found 355 pieces of
furniture across 47 sources; \textit{``Story continues below this ad''} appears in $75\%$
of one publisher's bodies. Republished
copies of one story are then grouped together, which I call near-duplicate grouping, since
a PTI report carried by three papers is one story. It reaches precision $0.86$ at recall
$0.81$ on 43 hand-judged pairs, and stops a syndicated copy of a training article from
leaking into validation. The splits hold 5{,}487 training, 1{,}120 validation and 1{,}159
test articles.
\section{Method}

Features are TF-IDF vectors, which score a term highly when it is frequent in one article
and rare across the corpus, so \textit{repo rate} counts for more than \textit{said}. The
base features are word n-grams, runs of one or two consecutive words.

I started with the taught models; logistic regression reached $0.752$ macro-F1 on the body
variant. K-nearest neighbours was excluded by reasoning rather than by a run: it is a lazy
learner comparing each query against all 5{,}487 training articles in roughly 30{,}000
dimensions, where the curse of dimensionality leaves distances nearly equal for every
pair. Tree ensembles need dense input, so the features were reduced to 256 components;
random forest reached $0.730$, extra trees $0.684$ and XGBoost $0.735$ at five times the
fit time.

The taught families plateaued near $0.75$, so I read further and tried a support vector
machine, which places the boundary so that the gap to the nearest training points of each
class is as wide as possible. That gap is the margin, and the score $w \cdot x + b$ is
the same linear quantity logistic regression pushes through the sigmoid. The plain SVM
only ties logistic regression, $0.753$ against $0.752$; the gain came from character
n-grams, runs of three to five consecutive characters, which lifted it to $0.771$ because
\textit{bowled}, \textit{bowler} and \textit{bowling} then share evidence. The final model
is a \texttt{LinearSVC} ($C=1$, \texttt{class\_weight="balanced"}) over word 1--2 and
\texttt{char\_wb} 3--5 grams, reading the title plus the first 4{,}000 body characters.

An SVM margin has no fixed scale and is not comparable across the 13 classes, so it must
never be shown as a confidence. Probability calibration adjusts scores so that a stated
confidence of $0.8$ is right about $80\%$ of the time. I use isotonic
regression~\cite{zadrozny}, a non-decreasing step function from raw score to observed
accuracy, over five grouped folds of train, and abstain below one global cut of $0.584$.
\section{Evaluation Protocol}

With $6.7\!:\!1$ imbalance, accuracy alone would hide failure on the small classes, so the
metrics come from the 13-class confusion matrix.
\begin{equation}
\text{Accuracy} = \frac{TP + TN}{TP + TN + FP + FN}
\end{equation}
\begin{equation}
P_i = \frac{TP_i}{TP_i + FP_i}, \qquad R_i = \frac{TP_i}{TP_i + FN_i}
\end{equation}
\begin{equation}
F1_i = \frac{2 \, P_i \, R_i}{P_i + R_i}
\end{equation}
\begin{equation}
\label{eq:macro}
\text{macro-F1} = \frac{1}{13} \sum_{i=1}^{13} F1_i
\end{equation}

Macro-F1 in (\ref{eq:macro}) weights every class equally, so education with 29 articles
counts as much as politics with 213. Weighted averaging lets the large classes dominate
and micro averaging reduces to accuracy, the trap this dataset sets.
The split is grouped and temporal: the cuts follow collection time and no near-duplicate
group crosses one. Intervals are bootstrapped over story groups rather than
articles~\cite{efron}, since copies of one wire story are not independent samples.
Comparisons use McNemar's test, a paired test over only the articles where two models
disagree~\cite{dietterich}. \textit{The Hindu} and \textit{The Guardian} are held out and
rescored on every run with the model refitted without them. The test split was opened
once, after the model was frozen.

ROC and AUC are not reported. AUC scores only how well a model ranks articles and does not
change if every probability shifts by a constant. This system acts on the probability
itself, comparing it against a cut of $0.584$, so a ranking score would miss the property
it depends on. I report expected calibration error instead, the gap between stated
confidence and observed accuracy.
\begin{table}[!tb]
\caption{Validation macro-F1: the body A/B test, then publisher holdouts.}
\label{tab:ladder}
\centering
\footnotesize
\setlength{\tabcolsep}{3.5pt}
\begin{tabular}{lcccc}
\toprule
Candidate & t+summary & t+body & Hindu & Guardian \\
\midrule
ComplementNB        & \textbf{0.635} & 0.594 & 0.589 & 0.563 \\
logistic regression & 0.698 & 0.752 & 0.747 & 0.736 \\
linear SVM, words   & 0.696 & 0.753 & 0.750 & 0.745 \\
\textbf{word+char SVM} & 0.712 & \textbf{0.771} & 0.769 & 0.724 \\
XGBoost, 256 comp.  & --- & 0.735 & 0.721 & 0.679 \\
random forest       & --- & 0.730 & 0.705 & 0.695 \\
MiniLM embeddings   & --- & 0.710 & 0.718 & 0.693 \\
\bottomrule
\end{tabular}
\end{table}

\section{Results}
\label{sec:results}

Table~\ref{tab:ladder} answers the main question. The body is worth $+0.059$ macro-F1 on
the final model, at McNemar $p = 7.5\times10^{-6}$ with disjoint intervals, and the gain
holds on every linear model. ComplementNB, a Naive Bayes variant adapted for imbalanced
text, is the one model that gets worse with more text, because its term-independence
assumption weakens as documents lengthen. Bodies were expected to carry publisher style
and hurt an unseen masthead; instead \textit{The Hindu} rose from $0.682$ to $0.769$ and
\textit{The Guardian} from $0.665$ to $0.724$.

No other family caught the linear model. Text is close to linearly separable in a very
high-dimensional sparse space, and reducing it to 256 dense components to make trees
usable discards the detail the linear model uses. MiniLM embeddings, a small neural
encoder that maps text to a fixed-length vector, landed at $0.710$, because it truncates
at 256 tokens and never reaches the body.

An oracle picking the correct member among the ten candidates would score $0.900$
$[0.878, 0.919]$, so the candidates do disagree usefully. Every majority vote still landed
on the incumbent: the best member set scored $0.772$ against $0.771$, at McNemar
$p = 1.0000$. Voting pays only when the members' errors are roughly independent, and here
they are asymmetric: MiniLM disagrees on 271 validation articles, rescues 70 and breaks
148 that were already right.

\begin{table}[!tb]
\caption{The held-out test split, opened once.}
\label{tab:test}
\centering
\footnotesize
\begin{tabular}{lcc}
\toprule
& validation & test \\
\midrule
macro-F1                     & 0.769 [0.738, 0.794] & \textbf{0.751 [0.719, 0.780]} \\
accuracy                     & 0.795 & 0.778 \\
expected calibration error   & 0.039 & \textbf{0.021} \\
coverage / accuracy at cut   & 0.761 / 0.879 & \textbf{0.807 / 0.834} \\
\bottomrule
\end{tabular}
\end{table}

Calibration costs almost nothing in accuracy, $0.769$ against a raw $0.771$, and the
reliability curve sits above the diagonal, so the model is under-confident. One global cut
was enough; per-class cuts scored $0.879$ against $0.891$ at matched coverage. Declining
the least certain $19.3\%$ of test articles raises accuracy on the rest to $83.4\%$, and
the validation-to-test gap of $-0.021$ sits inside the $\pm 0.05$ guard.

\section{Limitations}

\textit{society\_lifestyle} scores F1 $0.43$ on validation and $0.44$ on test, because the
class is a grab-bag of community, labour and lifestyle stories; it still reaches $91\%$
precision at a high cut, so abstention protects it. \textit{education} has only 29
validation articles, so its $0.72$ is noise. About $18.6\%$ of validation errors sit on
class pairs the annotators themselves disagreed on, so macro-F1 has a real ceiling below
$1.0$. The window is four days, the system is English-only and India-heavy, and the bundle
is $148.3$~MB.

\section{Conclusion}

Reading the article body instead of the headline and summary is worth $+0.059$ macro-F1
at $p = 7.5\times10^{-6}$. Logistic regression, random forests, boosted trees and
embeddings were measured on the same splits and none beat a linear support vector machine
over word and character n-grams. On unseen data the model scores $0.751$ macro-F1 and
files $80.7\%$ of articles at $83.4\%$ accuracy.
\scriptsize
\begin{thebibliography}{00}
\bibitem{iptc} IPTC, ``Media Topics NewsCodes,'' Int. Press Telecommunications Council, 2024.
\bibitem{joachims} T. Joachims, ``Text categorization with support vector machines: learning with many relevant features,'' in \textit{Proc. ECML}, 1998.
\bibitem{zadrozny} B. Zadrozny and C. Elkan, ``Transforming classifier scores into accurate multiclass probability estimates,'' in \textit{Proc. ACM SIGKDD}, 2002.
\bibitem{guo} C. Guo, G. Pleiss, Y. Sun, and K. Q. Weinberger, ``On calibration of modern neural networks,'' in \textit{Proc. ICML}, 2017.
\bibitem{chow} C. K. Chow, ``On optimum recognition error and reject tradeoff,'' \textit{IEEE Trans. Inf. Theory}, vol. 16, no. 1, 1970.
\bibitem{dietterich} T. G. Dietterich, ``Approximate statistical tests for comparing supervised classification learning algorithms,'' \textit{Neural Comput.}, vol. 10, no. 7, 1998.
\bibitem{efron} B. Efron and R. J. Tibshirani, \textit{An Introduction to the Bootstrap}. Chapman \& Hall, 1993.
\end{thebibliography}

\end{document}
